\documentclass[11pt,a4paper]{article}

\usepackage[utf8]{inputenc}
\usepackage[T1]{fontenc}
\usepackage{lmodern}
\usepackage{amsmath,amssymb,amsthm}
\usepackage{mathtools}
\usepackage{geometry}
\usepackage{hyperref}
\usepackage{xcolor}
\usepackage{booktabs}
\usepackage{enumitem}
\usepackage{microtype}
\usepackage{seqsplit}

\geometry{margin=2.8cm}
\setlength{\emergencystretch}{3em}
\sloppy
\newcommand{\ttbr}[1]{\texttt{\seqsplit{#1}}}
\newcommand{\Zp}{\mathbb{Z}_2}
\newcommand{\N}{\mathbb{N}}
\newcommand{\nuTwo}{\nu_2}
\hypersetup{
  colorlinks=true,
  linkcolor=blue!50!black,
  citecolor=blue!50!black,
  urlcolor=blue!50!black,
  pdftitle={Exact Exponent Cylinders, Precision Budgets, and Singular Coding in Accelerated Syracuse Dynamics},
  pdfauthor={Piero Borgatta}
}

\newtheorem{theorem}{Theorem}
\newtheorem{proposition}[theorem]{Proposition}
\newtheorem{lemma}[theorem]{Lemma}
\newtheorem{corollary}[theorem]{Corollary}
\theoremstyle{definition}
\newtheorem{definition}[theorem]{Definition}
\newtheorem{remark}[theorem]{Remark}

\title{Exact Exponent Cylinders, Precision Budgets,\\
  and Singular Coding in Accelerated Syracuse Dynamics\\
  \large Version 6}
\author{Piero Borgatta\thanks{Independent researcher. Correspondence:
\texttt{info@pieroborgatta.com}; ORCID 0009-0001-8025-2405.}}
\date{24 September 2026}

\begin{document}
\maketitle

\begin{abstract}
This paper continues the published v5 methodology
retrospective \cite{borgatta_v5}. It reports exact arithmetic results in the
subsequent Lean development. Every nonempty positive-exponent word is
characterized by one exact residue congruence, including the final odd
endpoint bit; repeated blocks inherit a modulus growing with their length.
A fixed expansive Syracuse exponent word
consumes a precisely quantified number of $2$-adic precision bits each
time it matches an integer orbit. Compatible changes of affine phantom
label preserve that decreasing counter; unrestricted changes can reset it.
An explicit natural-number lasso demonstrates the failure of a proposed
rank built from three short-word counters. A paper-level cancellation
family shows that a constant affine defect need not bound destination
precision. A separate theorem constructs,
at every finite precision, positive odd inputs approaching the $2$-adic
singular point whose accelerated outputs remain different. We give the
correct domain for infinite positive-exponent coding and record necessary
cycle and finite-barrier constraints. These results identify a concrete
certificate-search problem, but they do not prove strict descent for every
positive integer or exclude every nontrivial cycle. The paper distinguishes
Lean declarations, paper-level deductions, and open research targets.
\end{abstract}

\begingroup\small
\tableofcontents
\endgroup

\section{Scope and notation}

The v5 paper \cite{borgatta_v5} is a published retrospective of the
earlier program. This version concerns the later arithmetic
development. Let
\[
  S(n)=\frac{3n+1}{2^{a(n)}},\qquad
  a(n)=\nuTwo(3n+1)
\]
for a positive odd integer $n$. Every iterate of $S$ is again a positive
odd integer. For a finite word $w=(a_1,\ldots,a_\ell)$ of positive
exponents, put $A=\sum_i a_i$, $P=3^\ell$, and $Q=2^A$. Its affine
constant $C_w$ is defined by $C_{()}=0$ and
\[
  C_{w\mathbin{\|}(a)}=3C_w+2^{\sum w}.
\]
When an orbit starting at $n$ has exactly the first $\ell$ exponents
$w$, its endpoint $m=S^\ell(n)$ satisfies
\begin{equation}\label{eq:affine}
  Qm=Pn+C_w.
\end{equation}
This identity is formalized as
\ttbr{evalSyracuseWord\_mul\_pow\_sum\_eq\_affine\_of\_matches}
in \texttt{CollatzBridge.lean}. A finite word is \emph{expansive} when
$P>Q$. This refers to its affine slope, not to a claim that every
individual matched endpoint exceeds its source.

\section{Exact finite exponent cylinders}

For a nonempty word $w$ with every $a_i>0$, let $\mathcal C_w$ be the
set of natural starting values whose first $\ell$ accelerated Syracuse
exponents are exactly $w$. The extra final bit in the following theorem
distinguishes an odd endpoint from an integral but even one.

\begin{theorem}[Exact cylinder congruence]\label{thm:cylinder}
For every $n\in\N$,
\begin{equation}\label{eq:cylinder}
  n\in\mathcal C_w
  \quad\Longleftrightarrow\quad
  (3^\ell n+C_w)\bmod 2^{A+1}=2^A.
\end{equation}
If $k\ge1$ and $w^k$ denotes $k$ consecutive copies of $w$, then
\begin{equation}\label{eq:repeated-cylinder}
  n\in\mathcal C_{w^k}
  \quad\Longleftrightarrow\quad
  (3^{k\ell}n+C_{w^k})\bmod 2^{kA+1}=2^{kA}.
\end{equation}
\end{theorem}
\begin{proof}
For a matching word, Equation~\eqref{eq:affine} gives
$3^\ell n+C_w=2^A m$ with $m$ odd, hence the stated residue. Conversely,
the residue supplies an odd natural endpoint $m$ satisfying this affine
identity. Reverse induction through the positive exponent word, using
the parity of each nonempty suffix constant, recovers each intermediate
odd quotient and proves that every prescribed valuation is exact. Apply
the same statement to $w^k$, whose exponent sum is $kA$, for
Equation~\eqref{eq:repeated-cylinder}.
\end{proof}

The Lean declarations are
\ttbr{syracuseWordMatches\_iff\_exactResidue} and
\ttbr{repeatPhantomWord\_matches\_iff\_exactResidue} in
\texttt{ExactCylinders.lean}; both pass the full library build. The
final $+1$ matters: for $w=[1,2]$, $A=3$ and $C_w=5$, so $n=3$ makes
$9n+5=32$ divisible by $8$ but does not realize the word, whereas
$n=11$ gives $9n+5=104\equiv8\pmod{16}$ and does realize it.

\begin{corollary}[One residue class and its density; paper-level]
For each such $w$, $\mathcal C_w$ is one odd residue class modulo
$2^{A+1}$. Its natural density among positive integers is
$2^{-(A+1)}$, or $2^{-A}$ relative to positive odd integers. The
positive integer points of this class are dense \emph{within} its
$2$-adic clopen cylinder, whose Haar measure in $\Zp$ is
$2^{-(A+1)}$. For $w^k$, replace $A$ by $kA$.
\end{corollary}
\begin{proof}
The coefficient $3^\ell$ is invertible modulo $2^{A+1}$, so
Equation~\eqref{eq:cylinder} has exactly one residue class. Since
$C_w$ is odd and $2^A$ is even, that class is odd. Counting one class
per period gives the natural densities. Every $2$-adic subcylinder of
this class contains positive representatives of the same integer
congruence, giving density within the cylinder.
\end{proof}

The congruence theorem is formalized; uniqueness, density, and the
topological description in the corollary are deductions in this paper.
Ross's earlier Zenodo v3 states the corresponding one-class result
\cite{ross2026}; our contribution here is its exact Lean integration with
the precision and trace developments, not a claim of mathematical
priority.

\section{An exact precision budget for one word}

Fix a nonempty expansive word $w$. Write $D=P-Q>0$ and
\[
  H_w(n)=Dn+C_w>0,\qquad p_w(n)=\nuTwo(H_w(n)).
\]
The rational formal fixed point of the word is $-C_w/D$. Since $D$ is
odd, $p_w(n)$ is also the $2$-adic precision of the difference between
$n$ and that point.

\begin{theorem}[Matched-word precision tax]\label{thm:tax}
If $w$ matches the orbit at $n$ and has endpoint $m$, then
\begin{equation}\label{eq:tax}
  QH_w(m)=PH_w(n),\qquad p_w(m)+A=p_w(n).
\end{equation}
If $k$ consecutive copies of $w$ match, then
\begin{equation}\label{eq:budget}
  kA\le p_w(n)\le\lfloor\log_2 H_w(n)\rfloor.
\end{equation}
In particular, one fixed expansive word with $A>0$ cannot repeat
indefinitely along a positive integer orbit.
\end{theorem}
\begin{proof}
Multiply $H_w(m)=Dm+C_w$ by $Q$ and use
Equation~\eqref{eq:affine}:
\[
  Q(Dm+C_w)=D(Pn+C_w)+QC_w
    =P(Dn+C_w),
\]
because $D+Q=P$. The multiplier $P=3^\ell$ is odd, whereas $Q=2^A$.
Taking $2$-adic valuations gives Equation~\eqref{eq:tax}; iteration
telescopes, and nonnegativity of the final valuation gives
Equation~\eqref{eq:budget}. The final upper bound is the elementary
valuation-versus-bit-length inequality for a positive natural number.
\end{proof}

The Lean declarations are
\ttbr{syracuseWord\_precisionTax\_scaled\_identity\_of\_matches},
\ttbr{syracuseWordPhantomPrecision\_after\_eval\_add\_sum},
\ttbr{syracuseWordPhantomPrecision\_after\_iterate\_add\_mul\_sum}, and
\ttbr{no\_infinite\_matched\_expansive\_word} in
\texttt{PrecisionTax.lean}. \texttt{PrecisionTemplates.lean} supplies
the words $[1]$, $[1,2]$, and $[2,1]$, for which $H_w(n)$ is respectively
$n+1$, $n+5$, and $n+7$. The theorem is a bound on \emph{consecutive
copies of the same matched word}. It does not bound arbitrary switching
among different words.

\section{Switching labels and a concrete obstruction}

One can attempt to follow a changing formal point by assigning each state
an affine label $(D,C)$ and precision $\nuTwo(Dn+C)$. Suppose an exact
Syracuse step has exponent $e$ and sends $n$ to $m$:
\[
  2^e m=3n+1.
\]
If source and destination labels obey
\begin{equation}\label{eq:compat}
  3D'=qD,\qquad D'+2^e C'=qC
\end{equation}
for an odd factor $q$, then
\begin{equation}\label{eq:switch}
  2^e(D'm+C')=q(Dn+C).
\end{equation}
For nonzero numerators this implies exact precision loss by $e$.
This is \ttbr{switchingPhantomPrecision\_after\_step\_add\_exponent}
in \texttt{SwitchingPrecision.lean}.

The compatibility equations matter. At any odd $n$, a freely chosen
shift $C_k=(2^k-1)n$ gives $\nuTwo(n+C_k)=k$ for every $k$. Thus a
``choose the closest phantom afresh'' rule need not decrease. This is the
Lean theorem \ttbr{arbitrary\_phantom\_shift\_can\_set\_precision}.

\begin{proposition}[A short-trace lasso]\label{prop:lasso}
The exact matched segment $743\to1115\to1255$ has exponent blocks
$[1]$ and $[1,2]$. The three-counter vector
\[
  \bigl(\nuTwo(n+1),\nuTwo(n+5),\nuTwo(n+7)\bigr)
\]
takes the values
\[
  (3,2,1)\ \longrightarrow\ (2,5,1)\ \longrightarrow\ (3,2,1).
\]
The compatible hidden counter $\nuTwo(11n+19)$ loses four bits across
the full $[1,1,2]$ segment.
\end{proposition}

The segment and hidden loss are formalized by
\ttbr{concrete\_shortTrace\_lasso} and
\ttbr{hiddenTrace\_lasso\_tax} in
\texttt{ThreeTraceObstruction.lean}. The recurrence
$n_{k+1}=1024n_k+1767$ with $n_0=743$ produces arbitrarily precise
positive integer approximants to the formal point $-19/11$; it does not
produce an infinite positive integer orbit with a repeating state.
The repeated triple rules out a strictly decreasing rank determined
\emph{solely} by those three counters; it does not rule out an augmented
state or a rank tracking the compatible hidden counter.
Moreover, \texttt{AffineTraceGraph.lean} proves that two prescribed words
with different nondegenerate rational fixed points cannot both be
self-cycles of one scalar rational control state
(\ttbr{no\_common\_rational\_fixed\_point\_of\_distinct\_words}).
These facts constrain a certificate search based on a small fixed set of
phantom labels.

\section{The singular point and the coding domain}

\paragraph{Erratum to the published v4.}
The definition preceding the non-transport lemma in the published v4
\cite{borgatta_v4} called the accelerated Syracuse map a continuous map
$\widetilde S:\Zp\to\Zp$ and then called its valuation-word map a
topological conjugacy with a shift. The first assertion is false, so the
conjugacy argument as written is invalid. Its advertised global type is
also wrong: an even input can have first exponent $0$, outside the
positive-exponent alphabet, while at $x=-1/3$ the numerator vanishes
and has no finite valuation exponent. We replace those claims with the
finite arithmetic witness below and restrict infinite-word coding to the
invariant odd domain on which every numerator is nonzero. The logical
non-transport observation for positive integer orbits survives by the
direct arithmetic argument following the theorem; it requires no global
continuous extension or global topological conjugacy.

The total function \ttbr{Syracuse2adic} in \texttt{Syracuse2Adic.lean}
uses the odd unit factor of $3x+1$ and sets the value to $0$ at
$x=-1/3\in\Zp$. This total map agrees with $S$ on positive odd
integers, but it is not continuous at $-1/3$.

\begin{theorem}[Finite-precision separation]\label{thm:singular}
For every $R\in\N$, there are positive odd integers $n_1,n_7$ such that
\[
  2^R\mid(3n_1+1),\qquad 2^R\mid(3n_7+1),\qquad
  S(n_1)=1,\quad S(n_7)=7.
\]
\end{theorem}
\begin{proof}
Starting from $1$ and $9$, repeatedly apply $L(n)=4n+1$.
The identity $3L(n)+1=4(3n+1)$ raises the valuation by two without
changing the accelerated output. Thus, with
\[
  n_1(k)=\frac{4^{k+1}-1}{3},\qquad
  n_7(k)=\frac{28\cdot4^k-1}{3},
\]
both positive odd sequences have the fixed outputs $1$ and $7$, while
their numerators are divisible by $2^{2k+2}$. Choose $k$ large enough
for the given $R$.
\end{proof}

The Lean theorem is
\ttbr{singular\_outputs\_separated\_at\_every\_precision} in
\texttt{SyracuseSingularity.lean}. It asserts the finite arithmetic
separation, not a topological result. Since $3$ is a $2$-adic unit, both
sequences converge in $\Zp$ to $-1/3$. Their different constant outputs
then give the paper-level discontinuity conclusion.

Infinite positive-exponent words require a restricted domain. Let
$U=\{x\in\Zp:x\text{ odd},\ 3x+1\ne0\}$, and define
$D=\{x\in U:S^k(x)\in U\text{ for every }k\ge0\}$, with $S$ on $U$
given by the odd unit factor of $3x+1$.
For $a\ge1$ and odd $y\in\Zp$, the inverse branch
\[
  \phi_a(y)=\frac{2^a y-1}{3}
\]
is odd, has valuation exponent $a$, and maps to $y$. Compositions for a
fixed prefix of $W=(a_1,a_2,\ldots)$ agree modulo
$2^{a_1+\cdots+a_m}$. Their limit $\xi_W\in D$ is the unique point
with that exponent word. In particular, a positive odd integer equals
the point reconstructed from its own infinite word. The coding proof in
this paragraph is paper-level; the current Lean module formalizes the
arithmetic singularity witness, not this general reconstruction theorem.

Here \emph{not eventually periodic} means that no tail of a word repeats
with a fixed positive period. This qualification is essential: the word
of $n=3$ is $(1,4,2,2,\ldots)$, which is not periodic from the start but
is eventually periodic. For \emph{all} infinite positive-exponent words,
the statement ``$W$ is not eventually periodic $\Longrightarrow\xi_W$
is not a positive integer'' is equivalent to every positive odd integer
having a bounded accelerated orbit. Indeed, a bounded positive orbit is
pre-periodic and its exponent word is eventually periodic. Conversely,
suppose a tail repeats a nonempty block $w$ indefinitely. Then $P\ne Q$:
$P=3^{|w|}$ is odd and $Q=2^A$ is even. The expansive case $P>Q$ is
excluded for a positive integer by the matched-word precision budget
of Theorem~\ref{thm:tax}. In the remaining case $P<Q$, the block endpoint
recurrence $m=(Pn+C_w)/Q$ has slope strictly below one and bounded
iterates; its finitely many within-block intermediate states are bounded
as well. Thus the integer orbit is bounded. This is the
no-unbounded-orbit half of Collatz; the
coding equivalence does not prove it. A claim restricted to a particular
class of non-eventually-periodic words has different logical strength.

\section{Related work and priority}

Tao's almost-bounded-value theorem is an almost-all statement in
\emph{logarithmic} density \cite{tao2019}. Inselmann proves a different
almost-all descent estimate in \emph{natural} density on a logarithmic
time scale \cite{inselmann2024}. Neither is a pointwise theorem for every
positive integer. They set a stronger statistical context for the local
finite-precision identities here, while leaving the universal bridge open.

Sharpe's Lean repository reports near-periodicity and mechanical-itinerary
results, among other formalized Collatz statements \cite{sharpe2026}.
Fern\'andez and Ib\'a\~nez study extremal parity words through Christoffel
structures \cite{fernandez2026}. Ross's Zenodo v3 establishes that a
finite positive-exponent word occupies one odd residue class modulo
$2^{1+\sum a_i}$ \cite{ross2026}. These sources overlap with the word and
cycle language used here. The contribution claimed in this note is the
\emph{specific combination and formalization in this repository} of the
exact cylinder criterion, a matched-word precision budget, switching
compatibility, and explicit obstructions. We make no claim of absolute
mathematical priority
for the underlying coding or valuation observations.

Recent AI-assisted claims include Mazur's positive-density convergence
formalization hosted by ProofAtlas \cite{mazur2026} and Shaik's proposed
polylogarithmic descent for almost all starts \cite{shaik2026}.
ProofAtlas records a successful theorem-specific build while independent
publication review remains pending. Shaik reports a Palomar Registry v2
replay of four selected Lean declarations; that selected replay does not
by itself review the whole manuscript. The exact scope and evidence of
both claims merit separate audit. We use neither as a premise, and
neither presented claim gives a conclusion for every positive integer.

\section{Cycle and finite-barrier constraints}

If a nonempty word $w$ matches a positive integer $n$ and returns to
$n$, Equation~\eqref{eq:affine} gives
\begin{equation}\label{eq:cycle}
  (2^A-3^\ell)n=C_w.
\end{equation}
Consequently $2^A>3^\ell$, and any divisor of
$2^A-3^\ell$ divides $C_w$. These necessary statements are the Lean
theorems \ttbr{syracuseWordCycleEquationNat},
\ttbr{syracuseWordCycleContracting}, and
\ttbr{syracuseWordCycleModularSieve} in
\texttt{CycleConstraints.lean}. They do not show that all nontrivial
cycles are impossible.

For a matched contracting finite word, let $d=2^A-3^\ell>0$.
Rearranging Equation~\eqref{eq:affine} gives the exact displacement
identity
\[
  2^A(n-m)=dn-C_w.
\]
Hence $m<n$ if and only if $C_w<dn$. If the least nonnegative
representative $r=n\bmod 2^{A+1}$ satisfies $C_w<dr$, then $r\le n$
gives a finite residue certificate of strict descent. These are the
scope of \ttbr{evalSyracuseWord\_lt\_iff\_source\_above\_barrier}
and \ttbr{evalSyracuseWord\_lt\_of\_exactResidue\_above\_barrier}
in \texttt{FirstBarrier.lean}. Evaluating the threshold at the actual
source is algebraically equivalent to the desired drop; it is not an
independent global ranking principle.

The related module \texttt{A0InfiniteSubfamily.lean} proves descent on
the infinite residue classes $n=615+4096u$ and $n=2663+8192u$.
This is a genuine pointwise result on two families, with no assertion that
their union covers all positive odd integers.

\section{A testable continuation}

The results suggest a finite-certificate project with three separately
checkable obligations:
\begin{enumerate}[leftmargin=2em]
  \item Use Theorem~\ref{thm:cylinder} to enumerate finite exact exponent
        cylinders and attach each source residue either to a strict-descent
        certificate or to a declared control state.
  \item On every continuing edge, verify a compatible affine label
        transport of Equation~\eqref{eq:compat}, or explicitly account for
        any counter reset or cancellation. The lasso of Proposition~\ref{prop:lasso}
        invalidates a rank that records only the three shortest labels.
  \item Prove that every possible exit and every uncovered residue has an
        actual descent witness; a finite matrix bound or a count of
        exceptional residues is not that witness.
\end{enumerate}
The immediate falsifiable milestone is a small, exact certificate checker
that reports uncovered residue classes and incompatible label cycles,
not only successful finite examples. A theorem covering all positive odd
integers remains open. The named hypotheses
\ttbr{UniformStrictDescentHypothesis},
\ttbr{A0FirstBarrierExistsAll}, and
\ttbr{A0FirstBarrierThresholdAutomatic} have not been proved. The last
two together imply an exit for the A0 family only; the global descent
bridge still needs a cover of every positive odd source.

\paragraph{A proposed defect and cancellation ledger (paper-level).}
The compatible switching identity can be extended to measure a failure
of compatibility. On an exact edge $2^e m=3n+1$, assign integer affine
labels $F_i(n)=D_i n+C_i$ and $F_j(m)=D_j m+C_j$, and choose an odd
integer $q$. Define the edge defect
\[
  E(n)=(3D_j-qD_i)n+(D_j+2^e C_j-qC_i).
\]
Direct substitution gives
\[
  2^e F_j(m)=qF_i(n)+E(n).
\]
When the compatibility equations~\eqref{eq:compat} hold, $E$ is
identically zero and the verified switching identity applies. Otherwise
extend $\nuTwo$ to signed integers by
$\nuTwo(t)=\nuTwo(|t|)$ for $t\ne0$ and $\nuTwo(0)=+\infty$, and let
$a=\nuTwo(F_i(n))$, $b=\nuTwo(E(n))$ when both are nonzero. The exact
valuation ledger has three cases: if $a<b$, then
$\nuTwo(F_j(m))+e=a$; if $b<a$, then
$\nuTwo(F_j(m))+e=b$. If $a=b<\infty$, write
$qF_i(n)=2^a u$ and $E(n)=2^a v$ with $u,v$ odd, and put
$c=\nuTwo(u+v)\ge1$. Then
$\nuTwo(F_j(m))=a+c-e$, with the convention for a zero sum above.
The proposed scalar precision rank drops strictly on this edge only if
$c<e$; for $e=1$, this is impossible. Cancellation can make the
destination valuation arbitrarily higher, so the equal case must not
be treated as a bounded reset.

Indeed, take $e=1$, $q=3$, $F_i(n)=n+1$, and $F_j(m)=m+5$. Then
$E(n)=8$, a constant with no $2$-adic root. For every even $k\ge4$,
\[
  n_k=\frac{2^{k+1}-11}{3},\qquad
  m_k=S(n_k)=2^k-5
\]
are positive odd integers in the fixed exact $[1]$ cylinder
$n_k\equiv3\pmod4$. Direct
calculation gives
$\nuTwo(F_i(n_k))=\nuTwo(E(n_k))=3$ but
$\nuTwo(F_j(m_k))=k$. Thus absence of a defect root does \emph{not}
bound destination precision. This family is a paper-level arithmetic
counterexample, not a Lean declaration.

The defect still gives a useful exact cylinder test. On
$x\equiv r\pmod{2^K}$, write $E(x)=sx+b_0$ with $s\ne0$ and
$\alpha=\nuTwo(s)$. If $\nuTwo(E(r))<K+\alpha$, then
$\nuTwo(E(x))=\nuTwo(E(r))$ throughout the cylinder: write
$x=r+2^K t$ and compare the valuations of $E(r)$ and $s2^K t$.
Otherwise $E$
has a unique zero in that $2$-adic cylinder, approximated by positive
integer representatives with unbounded defect valuation. For $s=0$,
the defect is constant. A checker must also split the critical
$a=b$ cylinders, track zeros of the pulled-back destination numerator
$3D_j n+D_j+2^e C_j$, or supply a separate descent certificate. The
proposed falsifiable target is a finite exact-cylinder checker that
reports both types of valuation tower and candidate non-descending
control cycles at each tested resolution. No finite closure, global
cover, or termination theorem follows from this ledger alone.
The repository script \texttt{lean/scripts/phantom\_taxonomy/}
\texttt{switching\_defect\_probe.py} checks the identity, all three
valuation cases, and the constant-defect family at finite depth. It is
a reproducible diagnostic, not a proof of the proposed checker or a Lean
formalization.

\section{Verification boundary}

The cited declarations are in the \texttt{lean/CollatzShadowing} library.
The full library build completed on 24 September 2026 with Lean 4.29.1
and the matching Mathlib revision.
There are no source-level \texttt{sorry}, \texttt{admit}, or
user-declared \texttt{axiom} declarations in that library. Some finite
proofs use \texttt{native\_decide}, which places native compilation in
their trusted base; a source scan alone does not establish a kernel-only
axioms claim. A \texttt{\#print axioms} audit of selected structural
headline theorems reports only standard Lean axioms.

The pinned toolchain predates the official Lean 4.32.2 kernel soundness
fix \cite{lean4322} and further Lean 4.34.0 soundness fixes and checks
\cite{lean4340}. No defect is known to be exercised by this development.
A fresh build with Lean 4.34.0 and a compatible Mathlib version, followed
by another axioms audit, is a separate verification task; this manuscript
does not present it as completed.

\section{Conclusion}

The exact-cylinder theorem, finite-word precision tax, and compatible
switching identity provide
exact arithmetic progress measures. The trace lasso, incompatible cycle
labels, cancellation family, and singularity witness delimit where simple
globalizations fail. The cycle and barrier theorems supply necessary and
local tests.
Their joint value is a sharper specification for a future certificate
that must cover every positive odd orbit. Neither half of the Collatz
conjecture is settled here.

\begingroup\small
\begin{thebibliography}{99}
\bibitem{borgatta_v4}
P.~Borgatta, \emph{Phantom Orbit Shadowing for the Collatz Conjecture:
Verified Core, Two Barriers, and a Repositioning} (v4), Zenodo,
\href{https://doi.org/10.5281/zenodo.20544464}{10.5281/zenodo.20544464},
2026.
\bibitem{borgatta_v5}
P.~Borgatta, \emph{Cross-AI, Lean-Verified Mathematics: A Case Study on
the Collatz Conjecture} (v5), Zenodo, DOI
\href{https://doi.org/10.5281/zenodo.20554750}{10.5281/zenodo.20554750},
2026.
\bibitem{tao2019}
T.~Tao, \emph{Almost all orbits of the Collatz map attain almost bounded
values}, Forum of Mathematics, Pi 10 (2022), e12;
\url{https://arxiv.org/abs/1909.03562}.
\bibitem{inselmann2024}
M.~Inselmann, \emph{An approximation of the Collatz map and a lower bound
for the average total stopping time}, arXiv:2402.03276v3, 2024;
\url{https://arxiv.org/abs/2402.03276v3}.
\bibitem{sharpe2026}
M.~Sharpe, \emph{Collatz at the Critical Line}, Lean repository,
\url{https://github.com/msharpe248/collatz}, accessed 24 September 2026.
\bibitem{fernandez2026}
Fern\'andez and Ib\'a\~nez, \emph{Christoffel words as extremal structures
in Collatz dynamics}, arXiv:2607.24844, 2026;
\url{https://arxiv.org/abs/2607.24844}.
\bibitem{ross2026}
M.~Ross, \emph{Spike Structures and 2-adic Transition Laws in the
Accelerated Collatz Map}, Zenodo v3, 30 August 2026;
\href{https://doi.org/10.5281/zenodo.22181823}{10.5281/zenodo.22181823}.
\bibitem{mazur2026}
L.~Mazur, \emph{Explicit Positive-Density Collatz Convergence in
Logarithmic Time}, manuscript v2.1, 6 September 2026, ProofAtlas release;
\href{https://www.proofatlas.ai/formalizations/positive-density-log-time-collatz/}{formalization page}
(independent publication review pending).
\bibitem{shaik2026}
I.~A.~Shaik, \emph{Polylogarithmic Descent for Almost All Collatz Orbits in
Natural Density}, Zenodo v3.2.4,
\href{https://doi.org/10.5281/zenodo.22130385}{10.5281/zenodo.22130385},
2026; selected-declaration replay described at
\href{https://shaikidris.github.io/}{the author's evidence page}.
\bibitem{lean4322}
Lean Language Reference, \emph{Lean 4.32.2 release notes},
\url{https://lean-lang.org/doc/reference/latest/releases/v4.32.2/}, 2026.
\bibitem{lean4340}
Lean Language Reference, \emph{Lean 4.34.0 release notes},
\url{https://lean-lang.org/doc/reference/latest/releases/v4.34.0/}, 2026.
\end{thebibliography}
\endgroup

\end{document}
